Let $\sigma(n)=\sum_{d\mid n} d$ be the sum-of-divisors function for $n\in\mathbb{N}$.
Call $n$ \emph{$k$-perfect} if $\sigma(n)=k n$ for some integer $k\ge 1$.
Question: along any sequence of $k$-perfect numbers $n\to\infty$, must the corresponding multipliers satisfy
\[
k = o(\log\log n),
\]
meaning: for every $\varepsilon>0$ there exists $N(\varepsilon)$ such that if $n\ge N(\varepsilon)$ is $k$-perfect then $k\le \varepsilon\log\log n$.
Edge cases: $n=1$ is $1$-perfect; $k=2$ are the usual perfect numbers.
